\documentclass[10pt,a4paper]{amsart}
\usepackage[margin=19mm]{geometry}
\usepackage{amsmath,amssymb,mathtools}
\usepackage[scr=rsfs,cal=euler]{mathalfa}
\usepackage{xfrac}
\usepackage[unicode,hidelinks,bookmarks=false]{hyperref}
\newcommand{\ie}{i.e.\ }
\newcommand{\cf}{cf.\ }
\newcommand{\resp}{respectively\ }
\newcommand{\Subsection}{\subsection}
\providecommand{\nobreakdash}{\nobreak\hbox{-}}
\setlength{\emergencystretch}{2em}
\multlinegap0pt
\usepackage{mathtools}
\usepackage{amssymb}
\usepackage[shortlabels]{enumitem}
\usepackage[normalem]{ulem}
\usepackage{float}
\newtheorem{assumption}{Assumption}
\newtheorem{proposition}{Proposition}
\newtheorem{theorem}{Theorem}
\newcommand{\NN}{\mathcal{N}}
\title{Adaptive Partitioning and Learning for Stochastic Control of Diffusion Processes}
\author{Hanqing Jin, Renyuan Xu, Yanzhao Yang}
\date{Source: arXiv:2512.14991v3 (CC BY 4.0)}
\begin{document}
\maketitle

\noindent\emph{Source and licence.} Adaptive Partitioning and Learning for Stochastic Control of Diffusion Processes. Version arXiv:2512.14991v3 (CC BY 4.0), distributed by arXiv under the Creative Commons Attribution 4.0 International licence, http://creativecommons.org/licenses/by/4.0/, as recorded in the arXiv licence metadata for this exact version and shown on the abstract page https://arxiv.org/abs/2512.14991v3. Full licence notice: \texttt{licenses/arxiv-2512.14991-CC-BY-4.0.txt}.\par\smallskip

\noindent\textit{Editorial scope.} Counted unit: the article's theorem together (source0.tex lines 720--728). The article's complete original proof of it is delivered with it (source0.tex lines 2630--2751, one \texttt{proof} environment of 122 lines, which the article places away from the statement). No mathematics is written, repaired or re-derived: the statement and the proof are the article's own lines, in the article's own notation, and no retained line is substituted or re-flowed.  Retained uncounted as the source-local context the proof needs: lines 332--344; lines 346--357; lines 398--404; lines 676--693; lines 758--763; lines 808--814; lines 827--833; lines 888--897; lines 901--915; lines 959--967; lines 2539--2539; lines 2606--2606; lines 2627--2627.  Nothing was omitted from any retained span, so the package carries no omission record.  No retained line contains a citation, so the wrapper carries no bibliography and no bibliography entry.  No retained line contains a figure environment, an image-inclusion command or drawing code, so no image is delivered and the package declares no asset dependency.  Editorial formatting only, in addition: an AMS \texttt{amsart} preamble that restates the article's own declarations and macros used by the retained lines, namely the environments \texttt{assumption}, \texttt{proposition}, \texttt{theorem} on separate counters, together with the standard packages those lines need.  The retained environments print in this document as Assumption 1, Proposition 1, Theorem 1 in order of appearance, whereas the article prints the same items with the numbering of its own full document.  The complete native source of the article as distributed in the arXiv e-print for this exact version is bundled unchanged as \texttt{sources/191138/source0.tex}.
\begin{assumption}[Regularity of the dynamics]\label{ass:lipschitz}
Assume there exists constants $\ell_{\mu},\ell_{\sigma}>0$,  $m\in \mathbb{N}$ and $L_{0}> 0$ such that for all $h\in [H-1]$, $x_1,x_2 \in \mathbb{R}^{d_{\mathcal{S}}}$, and $a_1,a_2\in \mathcal{A}$, it holds that:
    \begin{eqnarray*}
        \|\mu_h(x_1,a_1)-\mu_h(x_2,a_2)\|&\leq&\ell_{\mu}\Big(\|x_1-x_2\|+\|a_1-a_2\|\Big),\nonumber\\
        \|\sigma_h(x_1,a_1)-\sigma_h(x_1,a_2)\|&\leq& \ell_{\sigma}\Big(\|x_1-x_2\|+\|a_1-a_2\|\Big),\nonumber\\
          \max_{h\in[H-1]}\{\|\mu_{h}(0,0)\|,\|\sigma_{h}(0,0)\|\}&\leq& L_{0}. 
    \end{eqnarray*} 

In addition, assume the following elliptic condition, i.e., there exists a  constant $\lambda>0$ such that $\forall x \in \mathbb{R}^{d_{\mathcal{S}}}, a \in \mathcal{A},$ and $  h \in [H-1]$:
\begin{eqnarray}\label{ass-eqn:volatility}
    \sigma_{h}(x,a)\sigma_{h}(x,a)^{\top}\succ \lambda I_{d_{\mathcal{S}}}.
\end{eqnarray}
\end{assumption}

\begin{assumption}[Regularity of the reward]\label{ass:expected reward local lipschitz}
Assume the expected reward is local Lipschitz, namely,  there exists  constants $\ell_{r}>0$, $m\in \mathbb{N}$ and $L_{0}> 0$ such that for all $h\in [H]$, $ x_1,x_2 \in \mathbb{R}^{d_{\mathcal{S}}}$, and $ a_1,a_2\in \mathcal{A}$, it holds that:
    \begin{eqnarray*}
    |\bar{R}_{h}(x_1,a_1)-\bar{R}_{h}(x_2,a_2)|&\leq&
{\ell_{r}\Big(\|x_1\|^m+\|x_2\|^m{+1}\Big)\,\Big(\|x_1-x_2\|+\|a_1-a_2\|\Big)},\nonumber\\
     \max_{h\in[H]}|\bar{R}_{h}(0,0)|&\leq&L_{0}.\label{ass-eqn:expected reward local lipschitz}
     \end{eqnarray*}
 In addition, assume that the reward distribution has sub-Gaussian tail decay, i.e., there exists a known constant $\theta>0$ such that $\forall x \in \mathbb{R}^{d_{\mathcal{S}}}, a \in \mathcal{A}, \lambda_{1} \in \mathbb{R}$, and $h \in [H]$:
 \begin{eqnarray}\label{ass-eqn:subGaussian reward}
    \mathbb{E}_{_{r_h\sim R_h(x,a)}}\Big[\exp\Big(\lambda_{1}(r_{h}-\bar{R}_{h}(x,a))\Big)\Big]\leq e^{\frac{\theta \lambda_{1}^2}{2}}.
\end{eqnarray}
\end{assumption}

\begin{proposition}[Local Lipschitz property of the value function]\label{thm:value function local Lipschitz}
Suppose Assumptions \ref{ass:lipschitz} and \ref{ass:expected reward local lipschitz} hold. Then for each  $h \in [H]$, it holds that
\begin{eqnarray}\label{eq:value function local Lipschitz}
   | V_{h}^{*}(x_1)- V_{h}^{*}(x_2)|\leq \overline{C}_{h}\Big(1 + \|x_1\|^m +\|x_2\|^m\Big)\|x_1-x_2\|,
\end{eqnarray}
 with $\overline{C}_{h}:=\overline{C}_{h}(\overline{C}_{h+1},L_{0},\ell_{\mu},\ell_{\sigma},\ell_{r},\Delta,m)$.
 \end{proposition}

\begin{theorem}\label{thm:high_prob_w_bound} Given Assumption \ref{ass:lipschitz}, it holds with probability at least $1-2\delta$ that, for any $ (h,k) \in [H-1]\times [K]$, $B \in \mathcal{P}_h^k$ with $n_h^k(B)>0$, and any $(x,a) \in B$,  

    \begin{eqnarray}
&&\mathcal{W}_{2}\Big(\NN(\widehat{\mu}_{h}^{k}(B)\Delta,\widehat{\Sigma}_{h}^{k}(B)\Delta),\NN(\mu_{h}(x,a)\Delta,\Sigma_{h}(x,a)\Delta)\Big) \nonumber\\
    &\leq& \Delta\kappa_{\mu}(\delta,\|\tilde{x}(^{o}B)\|,n_h^k(B))+ \frac{\Delta^{\frac{3}{2}}}{\sqrt{\lambda}}\kappa_{\mu}(\delta,\|\tilde{x}(^{o}B)\|,n_h^k(B))^{2}+\frac{\sqrt{d_{\mathcal{S}}}\Delta^{\frac{1}{2}}}{\sqrt{\lambda}}\kappa_{\Sigma}(\delta,\|\tilde{x}(^{o}B)\|,n_h^k(B))\nonumber\\
    &&\hspace{12mm}+
    \Big\|\overline{\mathbb{E}}[\widehat{\mu}_{h}^{k}(B)]-\mu_{h}(x,a)\Big\|\Delta+\Big\|\overline{\mathbb{E}}[\widetilde{\Sigma}_{h}^{k}(B)]-\Sigma_h(x,a)\Big\|\frac{\sqrt{\Delta}}{\sqrt{\lambda}},\label{eq:high_prob_w_bound}
   \end{eqnarray}

where  $\kappa_{\mu}:(0,1]\times (\mathbb{R}_+\cup\{0\})\times \mathbb{N}_+ \mapsto \mathbb{R}_+$ is defined as 
\begin{eqnarray*}
    \kappa_{\mu}(\delta,y,n) :=\frac{\eta(y)}{\sqrt{\Delta}}\Big(\sqrt{\frac{d_{\mathcal{S}}}{n}} + \sqrt{\frac{2\log(\frac{HK^2}{\delta})}{n}}\Big);
\end{eqnarray*}
$\kappa_{\Sigma}:(0,1]\times (\mathbb{R}_+\cup\{0\})\times \mathbb{N}_+ \mapsto \mathbb{R}_+$ is defined as 
\begin{eqnarray*}
    \kappa_{\Sigma}(\delta,y,n) :=\eta(y)^{2}\Bigg( D_1\Big( \sqrt{\frac{d_{\mathcal{S}}}{n}}+\frac{d_{\mathcal{S}}}{\sqrt{n}}\Big) +\Bigg( \sqrt{\frac{\log(\frac{D_2}{d_{\mathcal{S}}})}{D_3n}}+\frac{\log(\frac{D_2 HK^2}{\delta})}{D_3\sqrt{n}}\Bigg)\Bigg).
\end{eqnarray*}
\end{theorem}

\begin{theorem}\label{thm:reward high prob bound}
Under Assumption \ref{ass:expected reward local lipschitz}, It holds with probability at least $1-\delta$ that, for any  $(h,k) \in [H]\times [K], B\in \mathcal{P}_h^k$ with $n_h^k(B)>0$, and any $(x,a) \in B$, 
\begin{eqnarray}\label{eq:reward high prob bound}
    |\widehat{R}_{h}^{k}(B)-\bar{R}_h(x,a)| \leq \mbox{\rm R-UCB}_{h}^{k}(B)+\left|\frac{\sum_{i=1}^{n_h^k(B)}\bar{R}_{h}(X_{h}^{k_{i}},A_{h}^{k_{i}})}{n_h^k(B)}-\bar{R}_h(x,a)\right|.
\end{eqnarray}
\end{theorem}

\begin{theorem}\label{thm:BIAS Bound}
 With the same assumptions as in Theorem \ref{thm:high_prob_w_bound}, the following inequality holds for all $h\in [H-1],k\in [K]$, $B \in \mathcal{P}_h^k$ with $n_h^k(B)>0$, and any $(x,a) \in B$:
\begin{eqnarray}\label{eq:BIAS BOUND}
     \Big\|\overline{\mathbb{E}}[\widehat{\mu}_{h}^{k}(B)]-\mu_{h}(x,a)\Big\|\Delta+\Big\|\overline{\mathbb{E}}[\widetilde{\Sigma}_{h}^{k}(B)]-\Sigma_h(x,a)\Big\|\frac{\sqrt{\Delta}}{\sqrt{\lambda}} \leq \mbox{\rm T-BIAS}(B). 
\end{eqnarray}
\end{theorem}
The proof of Theorem \ref{thm:BIAS Bound} is deferred to Appendix \ref{app:BIAS Bound}.

\begin{theorem}\label{thm:RBIAS bound}
    With the same assumptions as in Theorem \ref{thm:reward high prob bound}, the following inequality holds for all $(h,k)\in [H]\times [K]$, $B \in \mathcal{P}_h^k$ with $n_{h}^{k}(B)>0$, and any $(x,a) \in B$:
    \begin{eqnarray}\label{eq:RBIAS bound}
\Big|\frac{\sum_{i=1}^{n_h^k(B)}\bar{R}_{h}(X_{h}^{k_{i}},A_{h}^{k_{i}})}{n_h^k(B)}-\bar{R}_h(x,a)\Big|\leq \mbox{\rm R-BIAS}(B). 
    \end{eqnarray}
\end{theorem}
The proof of Theorem \ref{thm:RBIAS bound} is deferred to Appendix \ref{app:RBIAS bound}.

For $k=0, h\in [H],B\in \mathcal{P}_{h}^{0}, S=\Gamma_{\mathcal{S}}(B)$, and $x\in \mathbb{R}^{d_{\mathcal{S}}}$, the function estimators are initialized with
\begin{eqnarray}\label{eq:initial value for estimation}
        \overline{Q}_h^0(B)&:=& \widetilde{C}_{h}(1+(\|\tilde{x}(^{o}B)\|+D)^{m+1}),\nonumber\\
    \widetilde{V}_h^0(S)&:=&\widetilde{C}_{h}(1+(\|\tilde{x}(S)\|+D)^{m+1}),\nonumber\\
        \overline{V}_h^0(x) &:=& \widetilde{C}_h(1+\|x\|^{m+1}).
\end{eqnarray}
In addition, for $(h,k)\in [H]\times ([K]\cup\{0\})$, we set $\overline{Q}_h^k(\bar{Z}^{\complement})$ as:
\begin{eqnarray}\label{eq:initialize outside Q}
     \overline{Q}_h^k(\bar{Z}^{\complement}):=-\widetilde{C}_{h}(1+\rho^{m+1}).
\end{eqnarray}

For $k\ge1$, $B\in\mathcal P_H^k$, $S\in\Gamma_{\mathcal S}(\mathcal P_H^k)$, and $x\in\mathcal S_1$, let us define
\begin{eqnarray}\label{eq:Q,V estimation for h=H}
     \overline{Q}_{H}^{k}(B) &:=&\left\{\begin{array}{lll}\widehat{R}_{H}^{k}(B)+\mbox{\rm R-UCB}_{H}^{k}(B)+\mbox{\rm R-BIAS}(B) && \mbox{ if } n_H^k(B)>0\\
     \overline{Q}_{H}^{0}(B) &&
       \mbox{ if } n_H^k(B)=0, 
       \end{array}\right. \nonumber\\ 
        \widetilde{V}_{H}^{k}(S)&:=&\min\Big\{\widetilde{V}_{H}^{k-1}(S), \max_{B \in \mathcal{P}_H^k,S \subset\Gamma_{\mathcal{S}}(B) }\overline{Q}_{H}^{k}(B)\Big\} ,\nonumber\\ 
         V_{H,k}^{\rm local}(x, S)&:=& \widetilde{V}_{H}^{k}(S)+C_{H}\Big(1+\|x\|^m+\|\tilde{x}(S)\|^m\Big)\|x-\tilde{x}(S)\|, \nonumber\\ 
        \overline{V}_{H}^{k}(x)&:=&\min_{S \in \Gamma_{\mathcal{S}}(\mathcal{P}_H^k)} V_{H,k}^{\rm local}(x, S).
\end{eqnarray}
For $x \in \mathbb{R}^{d_{\mathcal{S}}}\setminus\mathcal{S}_{1}$, we define 
\begin{eqnarray}\label{eq:Q,V estimation for h=H outside}
    \overline{V}_{H}^{k}(x):=\overline{V}_{H}^{k}\left(\frac{\rho}{\|x\|}x\right)+C_{H}(1+\|x\|^{m}+\rho^{m})\left\|\left(1-\frac{\rho}{\|x\|}\right)x\right\|. 
\end{eqnarray}
This extrapolation ensures that $\overline V_H^k$ is globally continuous.

\begin{theorem} \label{thm: True Upper Bounded By Estimates}
Under Assumptions \ref{ass:lipschitz}-\ref{ass:expected reward local lipschitz}, with probability at least $1-3\delta$, it holds that for all $(h,k) \in [H]\times[K]$, 
\begin{eqnarray}\label{eq:upperbound}
     \overline{Q}_{h}^{k}(B)&\geq& Q_{h}^{*}(x,a),\,\, \mbox{for all $B \in \mathcal{P}_h^k $ and  $ (x,a) \in B$}, \nonumber\\
     \widetilde{V}_{h}^{k}(S)&\geq& V_{h}^{*}(x), \,\,\mbox{for all $S \in \Gamma_{\mathcal{S}}(\mathcal{P}_h^k)$ and $x \in S$}, \nonumber\\
     \overline{V}_{h}^{k}(x)&\geq& V_{h}^{*}(x), \,\,\mbox{for all $x \in \mathbb{R}^{d_{\mathcal{S}}}$}.    
\end{eqnarray} 
\end{theorem}
The proof of Theorem \ref{thm: True Upper Bounded By Estimates} is deferred to Appendix \ref{app:proof-5-3}. \\

\subsection{Proof of Theorem \ref{thm:BIAS Bound}}\label{app:BIAS Bound}

\subsection{Proof of Theorem \ref{thm:RBIAS bound}}\label{app:RBIAS bound}

\subsection{Proof of Theorem \ref{thm: True Upper Bounded By Estimates}}\label{app:proof-5-3}

\begin{theorem}\label{thm:transition kernel wasserstein local lipschitz all together}
Assume Assumption \ref{ass:lipschitz} 
holds. With probability at least $1-2\delta$, we have that, for any $(h,k)  \in [H-1] \times [K]$, $B \in \mathcal{P}_h^k$ with $n_h^k(B)>0$, and $(x,a) \in B$:
\begin{eqnarray}
     &&\hspace{12mm}\left|\mathbb{E}_{X \sim \bar{T}_{h}^{k}(\cdot|B)}[V_{h+1}^{*}(X)]- \mathbb{E}_{Y \sim {T}_{h}(\cdot|x,a)}[V_{h+1}^{*}(Y)]\right|\label{eq:transition kernel wasserstein local lipschitz all together}\\
     &\leq&  \mbox{\rm T-UCB}_{h}^{k}(B)
     +L_{V}(\delta, \|\tilde{x}(^{o}B)\|)\,\,\Big(\Big\|\overline{\mathbb{E}}[\widehat{\mu}_{h}^{k}(B)]-\mu_{h}(x,a)\Big\|\Delta+\Big\|\overline{\mathbb{E}}[\widetilde{\Sigma}_{h}^{k}(B)]-\Sigma_h(x,a)\Big\|\frac{\sqrt{\Delta}}{\sqrt{\lambda}}\Big).\nonumber
\end{eqnarray}
\end{theorem}

\begin{proof} 
Recall from Theorems \ref{thm:transition kernel wasserstein local lipschitz all together} and \ref{thm:reward high prob bound}, we know that with probability at least $1-3\delta$,
\begin{eqnarray}\label{eq:simultanuous_result}
   \eqref{eq:transition kernel wasserstein local lipschitz all together}  \mbox{ and }   \eqref{eq:reward high prob bound} \mbox{ hold simultaneously}, 
\end{eqnarray}
 This fact serves as a building  block of the proof.\\
   
 %  We can now verify \eqref{eq:upperbound} by backward induction on $k$ (outer layer)   and forward induction on $h$ (inner layer). \\

For $k=0$, with the initialization of $(\overline{Q}_{h}^{0}$, $\widetilde{V}_h^0)_{h\in [H]}$ in  \eqref{eq:initial value for estimation}, we know \eqref{eq:upperbound} holds.
%for $k=0$ by the initialization of $\overline{Q}_{H}^{0}$, $\widetilde{V}_H^0$ defined in . 
Now assume that \eqref{eq:upperbound} holds for $k-1$ and we prove it holds for $k$.\\

\noindent \underline{For the case of $h=H$}: 
   For $B \in \mathcal{P}_H^k$ with $n_H^k(B)>0$ and for any $(x,a) \in B$, note that by \eqref{eq:reward high prob bound} and \eqref{eq:RBIAS bound} : 
    \begin{eqnarray}
          \widehat{R}_H^k(B)-\bar{R}_{H}(x,a) \geq -\mbox{\rm R-UCB}_H^k(B)-\mbox{\rm R-BIAS}(B). 
    \end{eqnarray}

    Therefore
$
          \overline{Q}_{H}^{k}(B)=\widehat{R}_{H}^{k}(B)+\mbox{\rm R-UCB}_{H}^{k}(B)+\mbox{\rm R-BIAS}(B)
          \geq Q_{H}^{*}(x,a).        
$    For $B\in \mathcal{P}_H^k$ with $n_H^k(B)=0$, by \eqref{eq:initial value for estimation}, 
                  we have $\overline{Q}_{H}^{k}(B)=\overline{Q}_{H}^{0}(B)\geq Q_{H}^{*}(x,a).$ So we proved the first inequality in \eqref{eq:upperbound}.

For any $S \in \Gamma_{\mathcal{S}}(\mathcal{P}_H^k)$ and any $x\in S$, 
%if $\widetilde{V}_{H}^{k}(S)=\widetilde{V}_{H}^{k-1}(S)$ then inequality \eqref{eq:upperbound} holds 
we have $\widetilde{V}_{H}^{k-1}(S)\ge V^*_H(x)$ by induction. 
%hypothesis on $k-1$, otherwise
Furthermore, 
   \begin{eqnarray*}
           \widetilde{V}_{H}^{k}(S)=\max_{B \in \mathcal{P}_H^k, \Gamma_{\mathcal{S}}(B) \supset S}\overline{Q}_{H}^{k}(B)\geq\overline{Q}_{H}^{k}(B^{*})\geq  Q_{H}^{*}(x,a_{H}^{*}(x))=V_{H}^{*}(x),
   \end{eqnarray*}
%in which, for $x\in S $, 
where $B^{*} \in \mathcal{P}_H^k$ is defined such that $(x,a_{H}^{*}(x)) \in B^{*}$. Hence 
we have $\widetilde{V}_{H}^{k}(S)\ge V^*_H(x)$. 

Finally, we check $\bar V^k_H(x)\ge V_{H}^*(x)$.  For any $x \in \mathcal{S}_1$, there exits  some $S^{\prime}\in \Gamma_{\mathcal{S}}(\mathcal{P}_H^k)$ such that
   \begin{eqnarray*}
           \overline{V}_{H}^{k}(x)&=&\widetilde{V}_{h}^{k}(S^{\prime})+C_{H}(1+\|x\|^m+\|\tilde{x}(S^{\prime})\|^m)\|x-\tilde{x}(S^{\prime})\|\\
           &\geq& V_{H}^{*}(\tilde{x}(S^{\prime}))+C_{H}(1+\|x\|^m+\|\tilde{x}(S^{\prime})\|^m)\|x-\tilde{x}(S^{\prime})\|\\
           &\geq& V_{H}^{*}(x),
   \end{eqnarray*}
where the last inequality holds since $|V_{H}^{*}(\tilde{x}(S^{\prime}))-V_{H}^{*}(x)|\leq C_{H}(1+\|x\|^m+\|\tilde{x}(S^{\prime})\|^m)\|x-\tilde{x}(S^{\prime})\|$ by \eqref{eq:value function local Lipschitz}.

For $x\in \mathbb{R}^{d_{\mathcal{S}}}\setminus \mathcal{S}_{1}$, by \eqref{eq:Q,V estimation for h=H outside}, we know that 
\begin{eqnarray*}
        \overline{V}_{H}^{k}(x)&=&\overline{V}_{H}^{k}\left(\frac{\rho}{\|x\|}x
        \right)+C_{H}(1+\|x\|^{m}+\rho^{m})\left\|\left(1-\frac{\rho}{\|x\|}\right)x\right\|\\
        &\geq& V_{H}^{*}(\frac{\rho}{\|x\|}x)+C_{H}(1+\|x\|^{m}+\rho^{m})\left\|(1-\frac{\rho}{\|x\|})x\right\|\\
        &\geq& V_{H}^{*}(x),
\end{eqnarray*}
where the first inequality holds since $\frac{\rho}{\|x\|}x\in \mathcal{S}_{1}$.\\

\noindent \underline{Induction ($h+1 \mapsto h$)}: 
Assume \eqref{eq:upperbound} holds for $h+1$ and we now show it also holds for $h$.


%When $k=0$, similar as before, \eqref{eq:upperbound} holds for $k=0$ by the initialization of $\overline{Q}_{h}^{0}$, $\widetilde{V}_h^0$ defined in \eqref{eq:initial value for estimation}. Assume that \eqref{eq:upperbound} holds for $k-1$, and let now prove it holds for $k$($k\geq 1$).

 
For $B\in \mathcal{P}_h^k$ with $n_h^k(B)>0$, by \eqref{eq:reward high prob bound} and \eqref{eq:RBIAS bound}, we have 
\begin{eqnarray}\label{eq:RUCB lower bound} 
          \widehat{R}_h^k(B)-\bar{R}_{h}(x,a) \geq -\mbox{\rm R-UCB}_h^k(B)-\mbox{\rm R-BIAS}(B). 
    \end{eqnarray}

We also have 
    \begin{eqnarray}\label{eq:TUCB lower bound}
     &&   \mathbb{E}_{X \sim \bar{T}_{h}^{k}(\cdot|B)}[\overline{V}_{h+1}^k(X)]- \mathbb{E}_{X \sim {T}_{h}(\cdot|x,a)}[V_{h+1}^{*}(X)]\nonumber\\
         &\geq& \mathbb{E}_{X \sim \bar{T}_{h}^{k}(\cdot|B)}[V_{h+1}^{*}(X)]- \mathbb{E}_{X \sim {T}_{h}(\cdot|x,a)}[V_{h+1}^{*}(X)]\nonumber\\
         &\geq& -\mbox{\rm T-UCB}_{h}^{k}(B)-L_{V}(\delta, \|\tilde{x}(^{o}B)\|)\mbox{\rm T-BIAS}(B). 
    \end{eqnarray}
The first inequality holds by induction hypothesis on $h+1$ and the second inequality holds due to \eqref{eq:transition kernel wasserstein local lipschitz all together} and \eqref{eq:BIAS BOUND}. 

Combining \eqref{eq:RUCB lower bound} and \eqref{eq:TUCB lower bound},  we have
%    \begin{eqnarray*}
  \begin{align*}
  \overline{Q}_{h}^{k}(B)
  ={}&\widehat{R}_{h}^{k}(B)+\mbox{\rm R-UCB}_{h}^{k}(B)
  +\mathbb{E}_{X \sim \bar{T}_{h}^{k}(\cdot|B)}[\overline{V}_{h+1}^{k}(X)]\\
  &+\mbox{\rm T-UCB}_{h}^{k}(B)+{\rm BIAS}(B)
  \geq Q_{h}^{*}(x,a).
  \end{align*}
 %   \end{eqnarray*}

For $B\in \mathcal{P}_h^k$ with $n_h^k(B)=0$, by \eqref{eq:initial value for estimation}, we also have
    \begin{eqnarray}
                  \overline{Q}_{h}^{k}(B)=\overline{Q}_{h}^{0}(B)\geq Q_{h}^{*}(x,a).
    \end{eqnarray}


   For any $S \in \Gamma_{\mathcal{S}}(\mathcal{P}_h^k)$ and any $x\in S$, 
   $\widetilde{V}_{h}^{k-1}(S)\ge V^*_h(x)$ holds by induction, and 
%   if $\widetilde{V}_{h}^{k}(S)=\widetilde{V}_{h}^{k-1}(S)$ then inequality \eqref{eq:upperbound} holds by induction hypothesis on $k-1$, otherwise
%   \begin{eqnarray*}
        %   \widetilde{V}_{h}^{k}(S)=
           $$\max_{B \in \mathcal{P}_h^k,\Gamma_{\mathcal{S}}(B) \supset S}\overline{Q}_{h}^{k}(B)\geq\overline{Q}_{h}^{k}(B^{*})\geq  Q_{h}^{*}(x,a_{h}^{*}(x))=V_{h}^{*}(x),$$
 %  \end{eqnarray*}
%in which for $x\in S $, 
where $B^{*} \in \mathcal{P}_h^k$ is the block containing $(x,a_{h}^{*}(x))$. 
%Recall that $a_h^*(x)$ is the optimal deterministic policy for state $x$ at timestamp $h$.
Hence $\widetilde{V}_{h}^{k}(S)\ge V^*_h(x).$

 

  Finally, for any $x\in \mathcal{S}_{1}$,  there exists a $S^{\prime}\in \Gamma_{\mathcal{S}}(\mathcal{P}_h^k)$ such that,
   \begin{eqnarray*}
           \overline{V}_{h}^{k}(x)&=&\widetilde{V}_{h}^{k}(S^{\prime})+C_h(1+\|x\|^m+\|\tilde{x}(S^{\prime})\|^m)\|x-\tilde{x}(S^{\prime})\|\\
           &\geq&V_{h}^{*}(\tilde{x}(S^{\prime}))+C_h(1+\|x\|^m+\|\tilde{x}(S^{\prime})\|^m)\|x-\tilde{x}(S^{\prime})\|\\
           &\geq& V_{h}^{*}(x),      
   \end{eqnarray*}
where the last inequality holds since $|V_{h}^{*}(\tilde{x}(S^{\prime}))-V_{h}^{*}(x)|\leq C_h(1+\|x\|^m+\|\tilde{x}(S^{\prime})\|^m)\|x-\tilde{x}(S^{\prime})\|$ by \eqref{eq:value function local Lipschitz}.

For $x\in \mathbb{R}^{d_{\mathcal{S}}}\setminus \mathcal{S}_{1}$, we know that 
\begin{eqnarray*}
        \overline{V}_{h}^{k}(x)&=&\overline{V}_{h}^{k}\left(\frac{\rho}{\|x\|}x\right)+C_h(1+\|x\|^{m}+\rho^{m})\left\|\left(1-\frac{\rho}{\|x\|}\right)x\right\|\\
        &\geq& V_{h}^{*}(\frac{\rho}{\|x\|}x)+C_h(1+\|x\|^{m}+\rho^{m})\left\|(1-\frac{\rho}{\|x\|})x\right\|\\
        &\geq& V_{h}^{*}(x), 
\end{eqnarray*}
where the first inequality holds since $\frac{\rho}{\|x\|}x\in \mathcal{S}_{1}$.
\end{proof}

\end{document}
